%@ pages 185-208
%@ status translated
% Original section 18, "Strictification", pp. 185--208.
%
% Translation notes (the mathematics is left exactly as printed):
%  * "foncteur large" is rendered "weak functor", as elsewhere in the
%    translation; "précat (de Segal)" is "precategory", as in section 2.
%  * "cmf" and "cmc" are both rendered "cmc".
%  * In the first paragraph the source prints $nSeCat$ (lower case) beside
%    $nSeCAT'$; both are kept as printed.
%  * In Lemma \ref{equivalence2} the source prints "Soit $L(M)$ un
%    remplacement fibrant ... pour $L(M)$", naming the replacement $L(M)$;
%    this is kept as printed (a % comment marks the spot).
%  * In part (2) of the proof of Theorem \ref{uneequivalence} the source
%    reads "ses composantes sont fibrantes et la flèche $02$ est une
%    cofibration"; "fibrantes" is kept as printed (a % comment marks the
%    spot).
%  * The source has "dans \cite{Hirschhorn} où Hirschhorn où prend soin",
%    with "où Hirschhorn" doubled; the duplication is not reproduced.

\clearpage
\section{Strictification}
\label{strictpage}

The question of the essential surjectivity of the morphism $\Phi$ of
Theorem \ref{correlation} poses the problem of ``strictifying'' a weak
functor $A: Y \rightarrow nSeCAT'$ into a strict functor $Y \rightarrow
nSeCat$. In the literature, the answers to this type of question are called
``strictification'' theorems, or ``coherence'' theorems.
See the discussion in \S 4 of Baez--Dolan \cite{BaezDolan};
there they refer to Gordon--Power--Street \cite{Gordon-Power-Street},
and to Mac Lane's coherence theorem \cite{MacLane2}.
See also Dunn \cite{Dunn97}.
%%%%%% added in v3: %%%%%%%%
See Cordier and Porter \cite{CordierPorter2} \cite{CordierPorter3}
who speak of ``rectification'' of diagrams, and Segal
(\cite{SegalTopology} Prop. B1).

A theorem very close to what we shall do in this chapter is
``A realization theorem 2.4'' of Dwyer--Kan \cite{DwyerKanDiags}.
In their theorem (which concerns simplicial presheaves, i.e.
Segal $0$-stacks), a ``weak functor'' is a functor
$Free _{\cdot}(Y)\rightarrow EnsSpl$ where $Free_{\cdot}(Y)$ is the standard
resolution by free categories. They show that every weak functor in this
sense is equivalent to a strict functor $Y\rightarrow EnsSpl$.

The theorem-type of this kind is the result of SGA 1 \cite{SGA1}
which says that a fibred category is equivalent to a split fibred category.
Our method of proof
for \ref{strictif3} will be
based on the method of SGA 1, but with a contribution from the techniques of
``Reedy categories'' coming from homotopy theory \cite{Reedy}
\cite{BousfieldKan} \cite{Hirschhorn} \cite{DHK} \cite{DwyerKanDiags}.
This method has already been used by the second author
to obtain a strictification result
for weak presheaves of topological spaces
(``flexible presheaves'') in \cite{flexible},
which turns out to be essentially equivalent to that of
\cite{DwyerKanDiags}.


One may envisage a slightly more general framework.
Let $M$ be a closed model category, and let $Y$ be a
$1$-category. Let $L(M)'$ be a fibrant replacement for the category
$L(M)$ (which is the Dwyer--Kan localisation of $M$ with respect to the weak
equivalences). Let
$F: Y \rightarrow L(M)'$ be a functor; one would like to find a functor
$G: Y \rightarrow M$ such
that
$p\circ G \sim F$ where $p: M\rightarrow L(M)'$ is the tautological morphism.
For the
closed model category $M=nSePC$, this would give the strictification for
functors $A:Y \rightarrow nSeCAT'$ in view of
Theorem \ref{intereqloc}.


We shall generalise further by considering, instead of a single
closed model category $M$, a left Quillen presheaf
$\mm$ over $Y ^o$. One recovers the result of the preceding paragraph by
taking the constant Quillen presheaf $\mm = \underline{M}$.

Recall that $\int '_YL(\mm )$ is a fibrant replacement for $\int
_YL(\mm )$ over $Y$
(i.e. the morphism to $Y$ becomes a fibrant morphism of
Segal $1$-categories).

\begin{definition}
\label{sestrictifient}
Let $Y$ be a small category and $\mm$ a left Quillen presheaf
on $Y^o$. We shall say that {\em the weak sections of $\mm$ over $Y$
strictify} if the morphism of Segal $1$-precategories
$$
Sect (Y, \int _Y\mm )\rightarrow Sect (Y, \int '_YL(\mm ))
$$
is essentially surjective.
\end{definition}



\begin{theorem}
\label{strictif1A}
If $\mm $ is a left Quillen presheaf on a Reedy category
$Y$ then the weak sections of $\mm$ over $Y$ strictify in the sense of
Definition \ref{sestrictifient}.
\end{theorem}

We shall give the proof below.

One should have the same result over an arbitrary category, provided
that the $\mm (y)$ admit sufficiently many limits. We have a
sketch of proof but there is a part that we did not want to carry out, so we
leave the statement in the form of a ``conjecture''.

\begin{conjecture}
\label{strictif1B}
Let $Y$ be a small category and $\mm$ a left Quillen presheaf
on $Y$. Suppose that $\mm$ satisfies hypothesis (o) which will be introduced
in \S 19 (which says in particular that the $\mm (y)$ admit arbitrary small
limits). Then the weak sections of $\mm$
over $Y$ strictify in the sense of Definition \ref{sestrictifient}.
\end{conjecture}

We shall propose a sketch of proof for this conjecture below.


\begin{corollary}
\label{strictif2A}
Let $Y$ be a Reedy $1$-category.
Let $M$ be a closed model category.
Then the morphism (of Segal $1$-categories)
$$
\underline{Hom} (Y,M)\rightarrow \underline{Hom}(Y,L(M)')
$$
is essentially surjective.
\end{corollary}
{\em Proof:} One applies Theorem \ref{strictif1A} to the constant Quillen
presheaf $\underline{M}$.
\eop

Conjecture \ref{strictif1B} would give the same corollary for every
category $Y$, provided one has hypothesis (o) of \S 19 for
$\underline{M}$.
In fact, we are interested in the statement in question for $M=nSePC$: the
following theorem will finally complete the proof of Theorem
\ref{correlation}.

\begin{theorem}
\label{strictif3}
Let $Y$ be a small $1$-category (not a Segal one). The weak sections of the
constant presheaf $\underline{nSePC}$ over $Y$ strictify
in the sense of Definition \ref{sestrictifient}.

In particular, every
object of $\underline{Hom}(Y, nSeCAT')$ is equivalent (in this
Segal $n+1$-category) to an object coming from a presheaf of
Segal $n$-categories on $Y$, i.e. the morphism $\Phi$ of Theorem
\ref{correlation} is essentially surjective.
\end{theorem}


{\em Proof of \ref{strictif3} using Conjecture
\ref{strictif1B}:}
One applies \ref{strictif1B} to the constant presheaf $\underline{M}$ for
the cmc
$M=nSePC$ of Segal $n$-precategories. It is easy to see that
$\underline{M}$ will satisfy
hypothesis (o). The tautological morphism
$$
M_f \rightarrow nSeCAT'
$$
sends equivalences to
equivalences, and hence extends (by the universal
property of the localisation) to a morphism
$$
L(M_f) \rightarrow nSeCAT'.
$$
The morphism $L(M_f)\rightarrow L(M)$ is a trivial cofibration,
as is the morphism to a fibrant replacement
$L(M)\rightarrow L(M)'$, and one obtains a
morphism
$$
L(M)' \rightarrow nSeCAT'.
$$
By Theorem \ref{intereqloc} this
morphism induces an equivalence
$$
L(M)' \stackrel{\cong}{\rightarrow} (nSeCAT')^{int, 1}.
$$
This morphism is by choice strictly compatible with the morphism from
$M_f$. The fact that $Y$ is a $1$-category (in particular
$1$-groupic) entails the equality
$$
\underline{Hom}(Y, nSeCAT')^{int, 1} =
\underline{Hom}(Y, (nSeCAT')^{int, 1}).
$$
An object $f\in \underline{Hom}(Y, nSeCAT')$ is therefore equivalent (in this
Segal $n+1$-category) to a morphism coming from
$g: Y\rightarrow L(M)'$. Applying the version of Corollary
\ref{strictif2A}
corresponding to \ref{strictif1B}, one obtains that $g$ is
equivalent (in $\underline{Hom}(Y, L(M)')$) to a morphism coming from
$h: Y\rightarrow M$; composing with a functorial fibrant replacement which
exists for $M=nSePC$, one obtains that $h$ can be chosen as a
morphism $h:
Y\rightarrow M_f$. The projection of $h$ in $\underline{Hom}(Y, nSeCAT')$ is
therefore equivalent to the projection of $g$, which in turn is equivalent
to $f$.

This statement is exactly the essential surjectivity of $\Phi$.
\eop

Since we can only propose a sketch of proof for \ref{strictif1B}, we shall
now give a complete proof of Theorem \ref{strictif3}.




\subnumero{Proof of Theorem \ref{strictif1A}}
We begin with the proof of Theorem \ref{strictif1A},
keeping the notation $\mm$ and $Y$. Our method of proof
consists in
a precise analysis of the diagrams indexed by a Reedy category,
using the ideas and techniques developed precisely
for this purpose by Reedy
\cite{Reedy}, Bousfield--Kan \cite{BousfieldKan}, Hirschhorn \cite{Hirschhorn},
Dwyer--Hirschhorn--Kan \cite{DHK}. One should regard our proof
as merely an application of the techniques of these references. We follow
more particularly the point of view of \cite{Hirschhorn}.

Part of the analysis has already been carried out in Lemma
\ref{reedydecomp} above, which we shall use now. Let $F^kY$ be the
filtration of $Y$ by
the full subcategories formed of the objects of degree $\leq k$. We
suppose given a section $\sigma$ of $\int '_YL(\mm )$ and we would like to
construct a section $u$ of $\int _Y\mm$ which is equivalent to $\sigma$.
We shall proceed by recursion on the degree,
and define $u$ on the
$F^kY$.
In the recursion we shall moreover choose $u$ fibrant and cofibrant for the
Reedy structure.

We suppose
therefore that we have found a suitable section $v$ on $F^kY$ and we want
to extend it to $u$ on $F^{k+1}Y$.
By \ref{reedydecomp}, $F^{k+1}Y$ decomposes (up to weak equivalence)
as the coproduct of $F^kY$
with categories of the form
$$
Y(y):= Latch (y) + \{ y\} + Match (y)
$$
for objects $y\in Y$ of degree $k+1$.
Set $F^kY(y):= F^kY\cap Y(y)$ whence
$$
F^kY(y)= Latch (y)  + Match (y).
$$
We shall extend $v|_{F^kY(y)}$ to a section $u_y$ on $Y(y)$,
equivalent to $\sigma |_{Y(y)}$. This will suffice to find $u$ and
the equivalence between $u$ and $\sigma |_{F^{k+1}Y}$ in view of
\ref{reedydecomp}.
On the other hand, the condition that $u$ be cofibrant and fibrant for the
Reedy structure is expressed in terms of the restrictions $u|_{Y(y)}$, i.e.
$u$ is cofibrant and fibrant if and only if (1) $v= u|_{F^kY}$ is, and
(2) for every $y$ of degree $k+1$, $u_y= u|_{Y(y)}$ is.

In sum, it suffices
to solve the problem of extending $v$ to $u$, for $F^kY(y) \subset
Y(y)$.
We have therefore reduced to the following situation: $Y=Y(y)$ and
$Z:= F^kY(y)$;
we have a section $\sigma$ of $\int '_YL(\mm )$ over $Y$, and a section
$v$ of $\int _Z\mm |_Z$ which is equivalent to $\sigma |_Z$; and we
seek a section $u$ over $Y$ which extends $v$ and which is equivalent
(as a section of $\int '_YL(\mm )$) to $\sigma$.
Recall the following notation:
$$
Y= Latch (y) + \{ y\} + Match (y);\;\;\;
Z= Latch (y) + Match (y).
$$

We have a morphism
$$
R: \int _{Latch(y)}L(\mm )\rightarrow L(\mm (y))
$$
corresponding to the
restriction morphisms $res_f: \mm (z)\rightarrow \mm (y)$
for every $f:z\rightarrow y$ in $Z$. This gives
(by extension along a trivial cofibration)
$$
R': \int '_{Latch(y)}L(\mm )\rightarrow L(\mm (y))' .
$$
Our section $\sigma |_{Latch (y)}$ gives a morphism
$$
R'\circ \sigma |_{Latch(y)} : Latch (y) \rightarrow L(\mm (y))'.
$$
The section $\sigma$ over $Latch(y)+ \{ y\}$ corresponds
(via a homotopy version, adapted to Segal categories, of the description of
Reedy diagrams of \cite{Hirschhorn} \cite{DHK}---a generalisation which
we use here without proof) to a morphism
$$
hocolim (R'\circ \sigma |_{Latch(y)})\rightarrow \sigma (y)
$$
in $L(\mm (y))'$. On the other hand, we also have the morphism
$$
{\bf r}: \int _{Latch(y)}\mm \rightarrow \mm (y)
$$
whence
$$
{\bf r}\circ v|_{Latch (y)}: Latch(y)\rightarrow \mm (y).
$$
The fact that $v|_{Latch(y)}$ is Reedy cofibrant, coupled with the condition
that the restriction functors be left Quillen functors,
implies that ${\bf r}\circ v|_{Latch (y)}$ is a $Latch(y)$-diagram of
$\mm (y)$ which is itself Reedy cofibrant. This implies that there exists a
natural equivalence in $L(\mm (y))'$:
$$
colim ({\bf r}\circ v|_{Latch (y)}) \cong
hocolim (R'\circ \sigma |_{Latch(y)})
$$
(in fact, this is the definition of $hocolim$, see \cite{BousfieldKan}
\cite{Hirschhorn} \cite{DHK}).

To have an extension of $v$ to a section $u_L$
over $Latch (y)+ \{ y\}$ it suffices (by the description of
\cite{Hirschhorn} \cite{DHK}) to have an object $u(y)\in \mm (y)$ and a
morphism
$$
colim ({\bf r}\circ v|_{Latch (y)})\rightarrow u(y).
$$
Moreover, $u_L$ is Reedy cofibrant (as is any extension of $u_L$ to $u$
over $Z$) if and only if the preceding morphism is a
cofibration. The section $u_L$ is equivalent
(as a section of the localised categories)
to $\sigma |_{Latch(y)+ \{ y\}}$ if there exists an equivalence in
$L(\mm (y))'$, $u(y)\cong \sigma (y)$, and a homotopy of commutativity for
the diagram
$$
\begin{array}{ccc}
colim ({\bf r}\circ v|_{Latch (y)})
& \cong & hocolim (R'\circ \sigma |_{Latch(y)})\\
\downarrow && \downarrow \\
u(y) & \cong & \sigma (y).
\end{array}
$$

Dually (although the argument is a little different) one studies the
extensions of
$v$ to a section $u_M$ over $Match (y)+ \{ y\}$. First, let us denote
$$
{\bf t}: \mm (y)\times Y \rightarrow \int _{Match (y)}\mm
$$
the morphism given by the restrictions. For
$$
y\stackrel{f}{\rightarrow}z\stackrel{g}{\rightarrow}w
$$
with composite $h:= gf$, one obtains, using the adjunctions, a
morphism (for $a\in \mm (z)$)
$$
res_f^{\ast} a \rightarrow res _h^{\ast} (res_ga).
$$
This gives a morphism
$$
{\bf t}^{\ast} : \int _{Match (y)}\mm \rightarrow
\mm (y).
$$
One likewise obtains on the localised categories
$$
T^{\ast} : \int '_{Match (y)}L(\mm )\rightarrow
L\mm (y)'.
$$
Composing with $\sigma$ one obtains
$$
T^{\ast}\circ \sigma |_{Match(y)}: Match(y)\rightarrow L(\mm (y)).
$$
The restriction of $\sigma$ to $\{ y\} + Match(y)$ corresponds
(again via
the homotopy version for Segal categories, of the description of
Reedy diagrams of \cite{Hirschhorn} \cite{DHK} which we do not prove
here) to a morphism
$$
\sigma (y) \rightarrow
holim (T^{\ast} \circ \sigma |_{Match(y)})
$$
in $L(\mm (y))'$.

Similarly, an extension $u_M$ of $v$ to $\{ y\} + Match(y)$
corresponds to an object $u(y)\in \mm (y)$ and a morphism
$$
u(y)\rightarrow lim ({\bf t}^{\ast} \circ v|_{Match(y)}).
$$
Here the limit exists because the indexing is finite.
The section $u_M$, as well as its possible extension $u$ to all of $Z$, is
Reedy fibrant if and only if this morphism is a fibration in $\mm
(y)$.

On the other hand, as above, one has an equivalence in $L(\mm (y))'$
$$
lim ({\bf t}^{\ast} \circ v|_{Match(y)})\cong
holim (T^{\ast} \circ \sigma |_{Match(y)})
$$
because $v|_{Match(y)}$ is Reedy fibrant by the recursion hypothesis, and the
adjoints $res _f^{\ast}$ are right Quillen functors, hence
${\bf t}^{\ast} \circ v|_{Match(y)}$ is a $Match(y)$-diagram of $\mm (y)$
which is
Reedy fibrant. An extension of the equivalence of $v$ with $\sigma |_Z$ to
an equivalence between $u_M$ and $\sigma |_{\{ y\} + Match (y)}$ corresponds
to the data of an equivalence $u(y)\cong \sigma (y)$ and of a homotopy of
commutativity for the diagram
$$
\begin{array}{ccc}
u(y) & \cong & \sigma (y) \\
\downarrow && \downarrow \\
lim ({\bf t}^{\ast} \circ v|_{Match(y)})& \cong &
holim (T^{\ast} \circ \sigma |_{Match(y)}).
\end{array}
$$

Let $x\in Latch(y)$ and $z\in Match(y)$
and denote by $f: x\rightarrow y$ and $g: y\rightarrow z$ the morphisms.
Then---
by the homotopy adjoint property (Lemma \ref{adjointh}) of the functors
induced by $res_g$ and $res_g^{\ast}$ between the simplicial categories $L(\mm
(y))$ and $L(\mm (z))$---the morphism
$$
res_{gf}(\sigma (x))\rightarrow \sigma (z)
$$
is equal to a morphism
$$
res_f(\sigma (x)) \rightarrow res_g^{\ast} (\sigma (z)).
$$
When $x$ varies in $Latch (y)$ and $z$ in $Match (y)$ one obtains the
following result.
Let us denote by $Latch(y)\sqcup Match(y)$ the disjoint union of these two
categories; it
is contained in $Z$ but in $Z$ there is, in addition, a (unique)
morphism from each object of $Latch(y)$ to each object of $Match(y)$.
An extension of $\sigma |_{Latch(y)\sqcup Match(y)}$ to $\sigma |_Z$
corresponds to a morphism
$$
hocolim (R'\circ \sigma |_{Latch(y)})\rightarrow
holim (T^{\ast} \circ \sigma |_{Match(y)}).
$$
If $\sigma |_Z$ is equivalent to the section induced by $v$, the map
above is induced by
$$
colim ({\bf r}\circ v|_{Latch (y)})\rightarrow
lim ({\bf t}^{\ast} \circ v|_{Match(y)}).
$$

The extension of $\sigma$ from $Z$ to $Y=Z\cup \{ y\}$ is determined
by a factorisation
$$
hocolim (R'\circ \sigma |_{Latch(y)})\rightarrow \sigma (y) \rightarrow
holim (T^{\ast} \circ \sigma |_{Match(y)}).
$$
This being given, in order to extend $v$ to a section $u$ on $Y$, one must
find a factorisation
$$
colim ({\bf r}\circ v|_{Latch (y)})\rightarrow u(y)\rightarrow
lim ({\bf t}^{\ast} \circ v|_{Match(y)})
$$
which induces the factorisation corresponding to $\sigma$.

We note moreover that $colim ({\bf r}\circ v|_{Latch (y)})$
is cofibrant (it is the colimit of a diagram which is cofibrant for the
Reedy structure, since ${\bf r}$ is a left Quillen functor);
and $lim ({\bf t}^{\ast} \circ v|_{Match(y)})$ is fibrant---it is the limit
of a diagram which is fibrant for the Reedy structure, since ${\bf t}^{\ast}$
is a right Quillen functor.

We have therefore reduced to the following problem: given a morphism
$$
h:a\rightarrow b
$$
in a closed model category $M=\mm (y)$, and a factorisation
$$
\overline{a} \rightarrow w \rightarrow \overline{b}
$$
in $L(M)$ where $\overline{a}$ and $\overline{b}$ are the images of $a,b\in M$
in $L(M)$, one wants to lift it to a factorisation
$$
a\rightarrow c \rightarrow b
$$
in $M$ with the first morphism cofibrant and the second
fibrant. One may suppose that $a$ is cofibrant and $b$ fibrant.

One may also suppose that one
has $w=\overline{c}$ in $L(M)$ with $c$ fibrant and
cofibrant. In this case, the morphisms of $L(M)$ come from morphisms
$f: a\rightarrow c$ and $g: c\rightarrow b$. One may suppose moreover that $g$
is a fibration. The factorisation in $L(M)$ is a homotopy between $gf$ and
$h$, as morphisms from $a$ to $b$. The homotopy may be
considered
as Quillen defined it, that is to say as a diagram
$$
i_0,i_1 :a\tworightarrows \alpha \stackrel{p}{\rightarrow} a ,\;\;\; k:\alpha
\rightarrow b
$$
with $pi_0=pi_1=1_a$ and $ki_0=gf$, $ki_1=h$ (and $p$ a trivial fibration, and
$i_0\sqcup i_1$ a cofibration). The morphism $i_0:a\rightarrow \alpha$
is a
trivial cofibration. The fact that we assumed
that $g$ is a fibration implies that there
exists a lift $w: \alpha \rightarrow c$ with $gw=k$
and $wi_0= f$. If one sets $f':= wi_1$, one obtains a homotopy (which we
still denote $w$) between $f$ and $f'$;
one also has $gf'=h$ and the homotopy $w$ composed with
$g$ is the homotopy $k$. In other words, the factorisation $h=gf'$ with
the identity as homotopy is equivalent to the factorisation which
corresponds to the homotopy $k$ between $gf$ and $h$. One obtains the
lift of our factorisation in $L(M)$ to a factorisation
$$
a\stackrel{f'}{\rightarrow}c\stackrel{g}{\rightarrow}b
$$
in $M$, which solves the problem.

To complete the proof of Theorem \ref{strictif1A},
let us apply the solution of the
problem above to the preceding situation with
$$
a= colim ({\bf r}\circ v|_{Latch (y)}),\;\;\;\;\;\;  b =
lim ({\bf t}^{\ast} \circ v|_{Match(y)});
$$
$$
\overline{a} = hocolim (R'\circ \sigma |_{Latch(y)}), \;\;\;\;\;\;
\overline{b}=  holim
(T^{\ast} \circ \sigma |_{Match(y)});
$$
and $w = \sigma (y)$. One finds the desired factorisation with $c=u(y)$.
\eop



\subnumero{The case of \texorpdfstring{$1$}{1}-stacks}

Before embarking on the sketch of proof for \ref{strictif1B}
and the proof of Theorem \ref{strictif3}, let us point out
the origin of
the idea by making the comparison with the case of $1$-stacks.

One can show that the $2$-category of $1$-stacks over $\Xx$
is equivalent (via Tamsamani's comparison \cite{Tamsamani} between the
$2$-categories following his definition, and Benabou's $2$-categories)
to the $2$-category $1CHAMP(\Xx )$ of $1$-stacks (not Segal ones) over
$\Xx$. If one adopts the definition of $1$-stack which uses the notion of
fibred category \cite{SGA1}, it is essentially a matter of showing that a
fibred category is equivalent to a split fibred
category, i.e. one which comes from a presheaf of categories.

This result is a theorem of SGA 1 \cite{SGA1}. The idea of the
proof is that if $\Ff \rightarrow Y$ is a fibred category
over a category $Y$, then the functor
$$
A:y\mapsto Sect^{eq} (Y/y, \Ff \times _Y(Y/y))
$$
(where $Sect^{eq}$ are the ``cartesian'' sections, i.e. those which send the
arrows downstairs to cartesian arrows upstairs) will be a strict functor from
$Y$ to
the $1$-category
$1Cat$ of $1$-categories, such that one has
an equivalence $\int _YA \cong \Ff $ of
fibred categories
over $Y$.

The key observation is that this strictification property comes
from the fact that the base $Y$ is a strict $1$-category.

This idea was taken up in \cite{flexible} to give a result of
strictification of ``flexible sheaves'' of topological spaces, which is
the analogue of Theorem \ref{strictif3} for the case $n=0$.

\subnumero{Idea for a direct proof of \ref{strictif3}}

We sketch here an idea for a direct proof of the strictification
for functors $\Xx ^o \rightarrow nSeCAT'$ (Theorem \ref{strictif3}),
following the case of $1$-stacks mentioned above.

Suppose now that $M= nSePC$ is the cmc
of Segal $n$-precategories. A functor $Y ^o\rightarrow M$ is a
Segal $n$-precategory over $Y$. On the other hand, one has \ref{intereqloc}
$$
L(M) \cong nSeCAT ^{int, 1}.
$$
Since $Y$ is already $1$-groupic, a functor
$Y ^o \rightarrow L(M) '$ is therefore, up to equivalence, the same thing
as a functor $Y ^o\rightarrow nSeCAT '$. Our question therefore becomes: is
every functor $Y ^o\rightarrow nSeCAT'$ equivalent to
a Segal $n$-precategory over $Y$?

The basic idea is to use the fact that $Y$ is a $1$-category which is
necessarily strict.
Let $F: Y ^o \rightarrow nSeCAT'$ be a morphism. For
$y\in Y$, we set
$$
G(y):= \Gamma (Y/y, F|_{Y/y}):= \underline{Hom}((Y/y)^o,
nSeCAT')_{1/}(\ast _y,  F|_{Y/y}).
$$
This varies functorially in $y$ and the values are fibrant Segal
$n$-categories. The problem is to show that $G$ is equivalent
to $F$.

We do not at present know how to find directly a morphism
between
$F$ and $G$ (and this is why this paragraph is only a sketch).
However, let us note for later that if $F$ already came
from a strict section, then there
would be an equivalence between $F$ and $G$ given
by the morphism
$$
F(y)\rightarrow \Gamma (Y/y, F|_{Y/y}).
$$

We shall circumvent the problem of finding a morphism between $F$ and $G$
by relying on Theorem \ref{strictif1A},
and by replacing $Y$ by a
Reedy category $B$ equipped with a subcategory $D\subset B$ such that
the localisation $L(B,D)$ is equivalent to $Y$. The idea that diagrams
over $Y$ are the same as diagrams over $B$ which are equivalences
in the direction of $D$ comes from the paper of Dwyer--Kan
\cite{DwyerKanDiags}.



\subnumero{Sketch of proof of Conjecture \ref{strictif1B}:}
Let $Y$ be a $1$-category and $\mm$ a left Quillen presheaf
on
$Y$ which satisfies conditions (o) of \S 19 below. We shall need the
following Reedy category. Let $\nu Y$ be the nerve of $Y$, and let
$B=\int _{\Delta} \nu Y$ be the category introduced for Lemma \ref{bd}.
Note that $B$ is also the category of simplices of $\nu Y$, denoted
$B=\Delta \nu Y$ in \cite{Hirschhorn} \cite{DHK}; by {\em loc. cit.}
$B$ is a Reedy category (the Reedy structure being induced by the
functor $B\rightarrow \Delta$). The morphism of simplicial sets
$$
\nu B\rightarrow \nu Y
$$
comes from a functor $B\rightarrow Y$, and moreover this functor sends
the morphisms of $D\subset B$ to the identities of $Y$. By Lemma
\ref{bd}, this induces an equivalence of simplicial or Segal
$1$-categories
$$
L(B,D)\cong 1SeL(B,D) \stackrel{\cong}{\rightarrow} Y.
$$

Consider the following diagram:
$$
\begin{array}{ccc}
Sect (Y, \int _Y\mm ) & \rightarrow & Sect (Y, \int '_YL(\mm ))\\
\downarrow && \downarrow \\
Sect (B, \int _B \varphi ^{\ast} \mm )&\rightarrow & Sect (B, \int '_B \varphi
^{\ast}L(\mm )).
\end{array}
$$
From the result \ref{strictif1A}, one obtains that the bottom morphism
is essentially surjective. On the other hand, the equivalence between $Y$
and $L(B,D)$ implies that the right-hand vertical morphism is fully
faithful with essential image the class
of sections $\sigma : B\rightarrow \varphi ^{\ast}L(\mm )'$ such that
$\sigma (g)$
is an equivalence for every morphism $g$ of $D$
(this idea comes from Dwyer--Kan \cite{DwyerKanDiags}). We shall call this
condition $D-eq$. If a section $\sigma $ is $D-eq$ and if $\sigma$ comes
from a section $u: B\rightarrow \int _B\varphi ^{\ast}\mm$ then $u$ is also
$D-eq$ in the sense that the restriction morphisms $res_g(u)$ are weak
equivalences
for $g\in D$. It suffices for us to prove that if
$u\in Sect (B, \int _B\varphi ^{\ast} \mm )$ is a $D-eq$ section in this
sense, then $u$ comes (up to weak equivalence for the structure of
model category) from a section $v\in Sect (Y, \int _Y\mm )$.

The functor
$$
p^{\ast}: Sect (Y, \int _Y\mm )\rightarrow Sect (B, \int _Bp^{\ast}\mm )
$$
admits a right adjoint
$$
p_{\ast}: Sect (B, \int _Bp^{\ast}\mm )
\rightarrow  Sect (Y, \int _Y\mm ).
$$
We claim that if $u\in Sect ^{D-eq}(B, \int _Bp^{\ast}\mm )$
is fibrant for the structure of type (II),
then one has the equivalence
$p^{\ast}p_{\ast} u\cong u$. For this we envisage a proof which would
follow the lines of the argument of \S 19 below (in particular this is
why one needs hypothesis (o), the essential point of which is the existence
of limits in the $\mm (y)$). This result would be in some sense a
relative version of the argument of \S 19. However, for the present version
of the paper
we have not verified the details of this argument and
this is why the
proof of \ref{strictif1B} contains a gap.

Granting this fact, one obtains that every
$u\in Sect ^{D-eq}(B, \int _Bp^{\ast}\mm )$ is equivalent to
a section of the form $p^{\ast} v$ for
$v\in Sect (Y, \int _Y\mm )$, which completes the proof in view of the
preceding remarks.
\eop





\subnumero{Proof of Theorem \ref{strictif3}}

Since there is a gap in our sketch of proof of
Conjecture \ref{strictif1B}, and since the idea for the direct proof given
earlier did not succeed either, we finally give a
complete proof by combining these two ideas. One had to apply
\ref{strictif1B} to the case of the constant Quillen presheaf $\mm =
\underline{M}$
with values $M=nSePC$, the model category of Segal $n$-precategories.
Note for example that $\int _Y\mm = Y\times M$ etc. With this notation, one
returns to the diagram
$$
\begin{array}{ccc}
Sect (Y, Y\times M ) & \rightarrow & Sect (Y, Y\times L(M )')\\
\downarrow && \downarrow \\
Sect (B, B\times M )&\rightarrow & Sect (B, B\times L(M )').
\end{array}
$$
Given $\sigma : Y\rightarrow L(M)'$,
one may apply the construction
of the paragraph before last, to obtain a section $G_{\sigma} :
Y\rightarrow
M$, essentially by $G_{\sigma}(y)= \lim _{Y/y} \sigma |_{Y/y}$. This is
compatible with the restriction to $B$ (via the morphism $\psi :
B\rightarrow Y$) in the sense that one has a morphism
$$
r:G_{\sigma} \circ  \psi \rightarrow G_{\sigma \circ \psi }.
$$
The fact that $\psi$ induces an equivalence $L(B,D)\cong Y$ implies
that for $b\in B$, one has
$$
\lim _{Y/\psi (b)}\sigma |_{Y/\psi (b)}\stackrel{\cong}{\rightarrow}
\lim _{B/b}(\sigma \circ \psi )|_{B/b},
$$
i.e. that the morphism $r$ is an equivalence (object-by-object over
$B$). On the other hand, by Theorem \ref{strictif1A}, $\sigma
\circ \psi $ is equivalent to a section coming from $u: B\rightarrow M$.
For strict sections, one has the desired morphism in the paragraph
before last. One therefore has an equivalence
$$
G_u \cong u,
$$
which gives
$$
G_{\sigma} \circ \psi \cong G_{\sigma \circ \psi} \cong \sigma \circ \psi .
$$
Setting $v:= G_{\sigma} : Y\rightarrow M$, one has solved (for this case) the
problem which gave rise to the gap in the proof of
Conjecture \ref{strictif1B}. One finishes the proof as before: the
fact that
$\psi$ induces an equivalence $L(B,D)\stackrel{\cong}{\rightarrow} Y$
implies that the morphism $Sect (Y,Y\times L(M)')\rightarrow Sect(B, B\times
L(M)')$ is fully faithful (cf \cite{DwyerKanDiags}).
Hence the equivalence $G_{\sigma}\circ \psi\cong \sigma \circ \psi$ comes
from an equivalence $G_{\sigma} \cong \sigma$.
This finally gives a complete proof of Theorem \ref{strictif3}.
\eop

Theorem \ref{strictif3} gives the essential surjectivity which was
missing
in order to complete the proof of Theorem \ref{correlation}.

{\em Remark--Exercise:}
\newline
In the proof of \ref{strictif3} one
could have
taken $B=\beta ^{\rm poset}(Y)$, see \ref{bdvariant}. In this case,
$B$ is the category underlying a partially ordered set;
because of this one would need Theorem \ref{strictif1A} only for this
type of category, which is a direct Reedy category. That
would simplify the proof of \ref{strictif1A} greatly, in particular
the matching objects $Match$ would be trivial and one would not need to
consider the connecting morphisms $Latch\rightarrow Match$. We leave to the
reader the exercise of writing a proof of \ref{strictif3}
(including the part of \ref{strictif1A} that would be needed) on the basis of
this remark.




\subnumero{An equivalence}

We now give an improvement of the preceding result
concerning
the essential surjectivity of the morphism \ref{sestrictifient}; this
morphism becomes an
equivalence if one localises (in the Dwyer--Kan sense) the model category
at the source of the morphism \ref{sestrictifient}. We give the statement
only for the case of
Reedy categories; we think that it remains true for an arbitrary base
category, under
hypotheses of type (o) of \S 19.

\begin{theorem}
\label{uneequivalence}
Let $Y$ be a Reedy category, and let $\mm$ be a left Quillen presheaf
on $Y$. Suppose that each $\mm (y)$ admits ``functorial
factorisations'' for the Quillen structure. Suppose the same thing for
$Sect (Z,\int _Z\mm |_Z)$ for every Reedy category $Z$ with a functor
to $Y$ (this functor not necessarily respecting the Reedy structure).
Then the
morphism
$$
L(Sect (Y,\int _Y\mm ))\rightarrow Sect (Y,\int '_YL(\mm ))
$$
is an equivalence.
\end{theorem}

\begin{corollary}
\label{pqcalclim}
Let $Y$ be a Reedy category and let $\mm$ be a left Quillen presheaf
on $Y$. Then one has an equivalence
$$
\lim _{\leftarrow , Y^o}L(\mm )\cong L(Sect ^{\rm eq}(Y,\int _Y\mm )).
$$
\end{corollary}
{\em Proof:}
In the equivalence of the theorem, a section $\sigma \in Sect (Y,\int _Y\mm
)$ goes to an $eq$-section of $\int 'L(\mm )$ if and only if the transition
morphisms $res_f\sigma (y)\rightarrow \sigma (z)$ (for $f:y\rightarrow z$
in $Y$) are weak equivalences; i.e. if and only if $\sigma$ is
an $eq$-section of $\int _Y\mm$. One therefore has (see \ref{stabilite}) an
equivalence:
$$
L(Sect ^{\rm eq}(Y,\int _Y\mm ))
\stackrel{\cong}{\rightarrow} Sect ^{\rm eq}(Y,\int '_YL(\mm )).
$$
On the other hand, one has by Proposition \ref{gammasect} that for the fibrant
replacement $L(\mm )'$ of $L(\mm )$,
$$
\Gamma (Y, L(\mm )')\rightarrow
Sect ^{\rm eq}(Y,\int '_YL(\mm )')
$$
is an equivalence; the morphism
$$
Sect ^{\rm eq}(Y,\int '_YL(\mm ))\rightarrow Sect ^{\rm eq}(Y,\int '_YL(\mm )')
$$
is an equivalence; and by Proposition \ref{calclim},
$$
\lim _{\leftarrow , Y^o}L(\mm ) \cong \Gamma (Y, L(\mm )').
$$
Hence the statement.
\eop

Before giving the proof of Theorem \ref{uneequivalence}
in general, we shall treat a few particular cases. We isolate
these lemmas because they show well, by
examples, what happens in \ref{uneequivalence}.

For these lemmas, we shall use the technique of ``homotopy function
complexes''
developed in Hirschhorn \cite{Hirschhorn}, Dwyer--Hirschhorn--Kan
\cite{DHK} and which has its origins in the work on the localisation
\cite{DwyerKan3} as well as in the work of Reedy \cite{Reedy}.
We briefly recall what
this is about.
This recall partially overlaps that of \S 8.

Here $M$ is a closed model category, and we want
to compute the homotopy types of the simplicial $Hom$ sets in the
Dwyer--Kan localisation $L(M)$, i.e. for $x,y\in M$ we want to compute
$L(M)_{1/}(x,y)$. For this, one introduces (in the references above)
the notion of a {\em fibrant simplicial resolution} $y\rightarrow {\bf y}$.
Here ${\bf y}$ is a simplicial object of $M$, i.e. an object ${\bf y}\in
M^{\Delta ^o}$, with a morphism $c^{\ast} y\rightarrow {\bf y}$ where
$c^{\ast} y$ is the constant simplicial object with values $y$. One requires
that
for $p\in \Delta$, $y\rightarrow {\bf y}(p)$ be a weak equivalence;
and that ${\bf y}$ be fibrant for the Reedy closed model category
structure (i.e. what is called here ``of
type (III)'') on $M^{\Delta ^o}$.
In this case, the simplicial set
$$
M_{1/}(x, {\bf y}):= \; p \mapsto M_{1/}(x, {\bf y}(p))
$$
is a simplicial set
which is naturally equivalent to $L(M)_{1/}(x,y)$
\cite{DwyerKan3}.
\footnote{
Technically the
natural morphism here is a morphism
$$
M_{1/}(x, {\bf y})\rightarrow L^H(M)
$$
in the localisation {\em by hammocks}, see \cite{DwyerKan3} 4.4 and 7.2.
Note that the argument of \cite{DwyerKan3} 7.2, which they give for
the case of a pair of resolutions ${\bf x}$ ${\bf y}$ (cosimplicial and
simplicial respectively), works equally well to provide the morphism which
we are looking for here in the case of a single resolution. In what follows
we shall use the notation $L(M)$ but to be technically correct (i.e. to have
the morphisms we speak of) one should read $L^H(M)$. }
The notation $M_{1/}(a,b):= Hom _M(a,b)$
denotes the set of morphisms for the category $M$.


We note now that the simplicial set $M_{1/}(x, {\bf y})$ is
fibrant, i.e. Kan. Indeed, the Kan extension condition for
$M_{1/}(x, {\bf y})$ becomes a lifting condition for morphisms
from $x$ to an arrow of the form ${\bf y}(p)\rightarrow {\bf w}$ where
${\bf w}$ is a (recursive) fibre product of components of ${\bf y}$ (this
fibre product depends on which extension condition, i.e. which
``horn'', one is
looking at; we leave to the reader the task of writing down precisely the
expression for the
fibre product in question). The condition that ${\bf y}$ is fibrant for the
Reedy structure implies that these morphisms ${\bf y}(p)\rightarrow {\bf w}$
are fibrant; and the condition that all the ${\bf y}(p)$ are equivalent to
$y$ implies (in view of the form of the fibre product which corresponds to a
``horn'') that ${\bf y}(p)\rightarrow {\bf w}$ is a weak equivalence.
Hence (since $x$ is cofibrant) every morphism from $x$ to ${\bf w}$ lifts to
a morphism $x\rightarrow {\bf y}(p)$, which gives the Kan extension
condition.

The same argument shows that a cofibration $x\rightarrow x'$ induces a
Kan fibration
$$
M_{1/}(x', {\bf y})\rightarrow M_{1/}(x, {\bf y}).
$$


The first lemma concerns the direct product: it is the case of Theorem
\ref{uneequivalence} where the base category $Y$ is discrete, equal
to a set. Despite
appearances, this case is not totally obvious
because of the fact that an infinite product of simplicial sets may fail to
have the right homotopy type, if the factors are not fibrant (Kan). This fact
was noted in particular by Jardine in \cite{JardineBool} (who gives a
counterexample).

\begin{lemma}
\label{equivalence1}
Let $\{ M_i\}$ be a collection of closed model categories
(indexed by a set). We equip the product $\prod _i M_i$
with the product closed model category structure: the fibrations,
cofibrations and weak equivalences are the morphisms which are so with
respect to
each variable. Let $L(M_i)'$ be fibrant replacements (as
Segal categories) of the Dwyer--Kan localisations $L(M_i)$.
Then the natural morphism
$$
L(\prod _i M_i)\rightarrow \prod _i L(M_i)'
$$
(induced by its restriction to $\prod _i M_i$) is an equivalence.
\end{lemma}
{\em Proof:}
We note first that the set of objects is the same on both sides;
it is therefore a matter of proving that the morphism is (homotopically)
fully
faithful. Let $(y_i)$ and $(z_i)$ be two objects of the product $\prod _i M_i$.
For each $i$ one chooses a fibrant simplicial resolution $z_i
\rightarrow
{\bf z}_i$. One thus obtains a fibrant simplicial resolution
$$
(z_i)\rightarrow ({\bf z}_i)
$$
for the product. Hence
$$
L(\prod _i M_i) _{1/}((y_i), (z_i))\cong
(\prod _i M_i)_{1/}((y_i), ({\bf z}_i)) = \prod _i \left(
(M_i)_{1/}(y_i,{\bf z}_i)\right) .
$$
On the other hand, for each $i$ the morphism
$$
L(M_i)_{1/}(y_i, z_i)\rightarrow L(M_i)'_{1/}(y_i,z_i)
$$
is an equivalence of simplicial sets with the second term fibrant.
One has for each $i$
$$
L(M_i)_{1/}(y_i, z_i)\cong (M_i)_{1/}(y_i,{\bf z}_i),
$$
and as we remarked above, the
$(M_i)_{1/}(y_i,{\bf z}_i)$ are fibrant. The direct product of fibrant
simplicial sets preserves homotopy equivalences
\cite{JardineBool}, hence
$$
\prod _i \left(L(M_i)'_{1/}(y_i, z_i)\right) \cong \prod _i \left(
(M_i)_{1/}(y_i,{\bf z}_i)\right) .
$$
Note that one has
$$
\prod _i \left(L(M_i)'_{1/}(y_i, z_i)\right) =
(\prod _i L(M_i)')_{1/}(y_i, z_i);
$$
one obtains the equivalence
$$
(\prod _i L(M_i)')_{1/}(y_i, z_i)\cong
L(\prod _i M_i) _{1/}((y_i), (z_i)).
$$
This equivalence is homotopic to the morphism
$$
L(\prod _i M_i) _{1/}((y_i), (z_i))\rightarrow
(\prod _i L(M_i)')_{1/}(y_i, z_i)
$$
of the lemma, which is (by definition) the one induced by the morphism
$$
(\prod _i M_i) _{1/}((y_i), (z_i))\rightarrow
(\prod _i L(M_i)')_{1/}(y_i, z_i),
$$
which sends weak equivalences to equivalences.
\eop

\begin{lemma}
\label{equivalence2}
Let $M$ be a closed model category. Let $I$ be the category
with objects $0,1$ and one morphism $0\rightarrow 1$, and let $M^I$ be the
category of the $I$-diagrams in $M$ (which admits a Reedy cmc structure, for
example). Let $L(M)$ be a fibrant replacement (as a
Segal category) for $L(M)$. % The source prints $L(M)$ twice here; kept as printed.
Then the natural morphism
$$
L(M^I)\rightarrow \underline{Hom}(I, L(M)')
$$
(induced by its restriction to $M^I$) is an equivalence. Here $\underline{Hom}$
is the internal one for Segal $1$-precategories.
\end{lemma}
{\em Proof:}
An arrow of $L(M)'$ is a homotopy class of morphisms of $M$,
hence every object of $\underline{Hom}(I, L(M)')$ is equivalent to an object
coming from $M^I$. It is therefore a matter of proving that the morphism in
question is
homotopically fully faithful. Let $y$ and $z$ be $I$-diagrams of
$M$, i.e. $y= (f: y_0\rightarrow y_1)$ with $y_i\in M$ and
$z= (g: z_0\rightarrow z_1)$.
We suppose that $y$ is cofibrant for that one
of the two possible Reedy structures for which
$f$ is a cofibration.
If ${\bf z}$ is a fibrant simplicial
resolution of $z$ then one has
$$
L(M^I)_{1/}(y,z) \cong M^I_{1/}(y, {\bf z}).
$$
One may write ${\bf z} = ({\bf g}: {\bf z}_0\rightarrow {\bf z}_1)$.


On the other hand,
$$
\underline{Hom}(I, L(M)')_{1/}(y,z)
$$
represents the functor of simplicial sets
$$
K\mapsto Hom ^{y,z}(I\times \Upsilon (K), L(M)')
$$
where the right-hand term is the subset of morphisms
$r:I\times \Upsilon (K)\rightarrow L(M)'$ with $r|_{I\times \{ 0\} } = y$
and $r|_{I\times \{ 1\} } = z$. A morphism
$$
I\times \Upsilon (K)\rightarrow L(M)'
$$
is equivalent (by dividing the square into two triangles---the reader
is invited to
draw a square with $f: y_0\rightarrow y_1$ on one side and $g:
z_0\rightarrow z_1$ on the other and with arrows labelled $K$ between $y_0$
and $z_0$
and between $y_1$ and $z_1$; divided into two triangles by an arrow labelled
$K$
from $y_0$ to $z_1$) to two morphisms
$$
\Upsilon ^2(\ast , K)\rightarrow
L(M)' $$ and
$$
\Upsilon ^2(K,\ast )\rightarrow L(M)'
$$
with the same restriction to the diagonal $\Upsilon (K)$. The
restriction morphisms of the simplicial sets of these diagrams, over
the simplicial set of diagrams for the diagonal $\Upsilon (K)$, are
fibrations. One obtains a formula with a homotopy fibre product,
i.e. a homotopy cartesian square
$$
\begin{array}{ccc}
\underline{Hom}(I, L(M)')_{1/}(y,z)&\rightarrow & L(M)_{1/}(y_0, z_0)\\
\downarrow && \downarrow \\
L(M)_{1/}(y_1, z_1) & \rightarrow &  L(M)_{1/}(y_0, z_1).
\end{array}
$$
There are two possible Reedy structures on $M^I$ (one may
choose $deg(0)>deg(1)$ or the opposite); we shall take the structure ($deg(0) <
deg (1)$) for which the cofibrant objects are the cofibrations
$y_0\rightarrow y_1$ between cofibrant $y_i$. In this case, and for a
fibrant simplicial resolution ${\bf z}$ (whose component ${\bf z}_1$
is in
particular also a fibrant simplicial resolution), the morphism of
composition with $f$
$$
M_{1/}(y_1, {\bf z}_1)\rightarrow M_{1/}(y_0, {\bf z}_1)
$$
is a fibration of simplicial sets (cf the discussion before these lemmas).
It follows that the cartesian square of simplicial sets
$$
\begin{array}{ccc}
M^I_{1/}(y,{\bf z}) & \rightarrow &  M_{1/}(y_0,{\bf z}_0)\\
\downarrow && \downarrow \\
M_{1/}(y_1,{\bf z}_1)& \rightarrow & M_{1/}(y_0,{\bf z}_1)
\end{array}
$$
is also homotopy cartesian. The three corners
of the bottom and of the right coincide
up to homotopy with the three corresponding corners in the square
above for the localised categories. One deduces the equivalence
$$
M^I_{1/}(y,{\bf z})\cong \underline{Hom}(I, L(M)')_{1/}(y,z)
$$
which gives the desired statement.
\eop





\subnumero{Proof of Theorem \ref{uneequivalence}}
Note that the morphism in question is essentially surjective by
\ref{strictif1A}. It is therefore a matter of proving that it is fully
faithful.

For the moment we shall suppose that the theorem is known for
{\em direct} or {\em inverse} categories, i.e.
categories with a function
``degree'' where all arrows except the identity are strictly
monotone for the degree. This will be justified later (cf point (1)
below). In fact, we shall make a first pass of our proof to
settle this case, then the second pass which we describe
now---the
reason for this contortion being that for the direct case we need
exactly the same arguments as for the general case, with even
a few simplifications but not enough to justify, for us who are
lazy, repeating the argument twice; so, one begins with the
most complicated version (this manoeuvre will be better explained in (1) at
the end of the proof). We discovered this technique
which consists in considering first the direct case and
then the Reedy case in \cite{Hirschhorn} where
Hirschhorn takes care to treat the ``direct'' case independently
first.

We shall proceed by recursion on the length of the degree function on $Y$.
We suppose that all objects of $Y$ are of degree $\leq k$, and that the
theorem is proved for Reedy categories with degree
function of length $\leq k-1$. We set $Z:= F^{k-1}Y$; in particular this
hypothesis applies to $Z$.

We shall denote by $y_i$ the objects of $Y$ of degree
$k$ (the index $i$ is in an index set which we do not put into the
notation). For each of
these objects, we denote
$$
Y_i:= Latch(y_i)+\{ y _i\} +Match (y_i),
$$
and $Z_i := Y_i \cap Z= Latch(y_i)+Match(y_i)$. Note that $Y_i$ and $Z_i$ are
direct (or else inverse) categories, so one may suppose that the
theorem applies to them.

Let $[0,1,2]$ be the category associated to
the ordered set $0 < 1 < 2$ and
$[0,2]$ its full subcategory with objects $0$ and $2$.

Let $\mm (y_i) ^{[0,1,2]}$
(resp. $\mm (y_i) ^{[0,2]}$)
be the category of $[0,1,2]$-diagrams (resp. $[0,2]$-diagrams)
in $\mm (y_i)$. We shall denote the diagrams $u=(u_0\rightarrow u_1\rightarrow
u_2)$ (or $u=(u_0\rightarrow u_2)$).

One has the following cartesian square of categories:
$$
\begin{array}{ccc}
Sect (Y_i,\int _{Y_i}\mm ) & \rightarrow & \mm (y_i) ^{[0,1,2]}\\
\downarrow && \downarrow \\
Sect (Z_i,\int _{Z_i}\mm ) & \rightarrow & \mm (y_i) ^{[0,2]}.
\end{array}
$$
The top morphism sends a section $\sigma$ to the diagram
$$
latch(\sigma ,y_i) \rightarrow \sigma (y_i) \rightarrow match (\sigma , y_i).
$$
The bottom morphism sends a section $\tau$ over $Z$ to the diagram
$$
latch(\tau  ,y_i) \rightarrow  match (\tau , y_i).
$$

We shall denote by $Sect (Y_i,\int _{Y_i}\mm )_{c,f}$ the subcategory of
sections
whose direct part (i.e. the restriction of the section to $Latch(y_i)+\{
y_i\}$) is cofibrant, and the inverse part (i.e. the restriction of the
section
to $\{ y_i\} +Match(y_i)$) is fibrant. Likewise for
$Sect (Z_i,\int _{Z_i}\mm )_{c,f}$. We shall denote by $\mm (y_i) ^{[0,1,2]}_{c,f}$
the category of diagrams $u$ in which
$u_0$ is cofibrant, $u_0\rightarrow u_1$ is
a cofibration, $u_2$ is fibrant and $u_1\rightarrow u_2$
is a fibration. We
shall denote by $\mm (y_i) ^{[0,2]}_{c,f}$ the category of diagrams $v$ with
$v_0$
cofibrant and $v_2$ fibrant.

The square above also gives a cartesian square
$$
\begin{array}{ccc}
Sect (Y_i,\int _{Y_i}\mm ) _{c,f}& \rightarrow & \mm (y_i) ^{[0,1,2]} _{c,f}\\
\downarrow && \downarrow \\
Sect (Z_i,\int _{Z_i}\mm )  _{c,f}& \rightarrow & \mm (y_i) ^{[0,2]} _{c,f}.
\end{array}
$$
For each of these categories one has an obvious notion of weak
equivalence---which we shall dispense with putting into the notation of the
localised categories. The advantage
of this second square is that the horizontal morphisms respect
weak equivalences (which is not the case for the first square,
since the
horizontal morphisms involve limits and colimits).

Now one also has the cartesian square of categories
$$
\begin{array}{ccc}
Sect (Y,\int _{Y}\mm )& \rightarrow & \prod _iSect (Y_i,\int _{Y_i}\mm )
\\ \downarrow && \downarrow \\
Sect (Z,\int _{Z}\mm ) & \rightarrow &
\prod _i Sect (Z_i,\int _{Z_i}\mm ).
\end{array}
$$
If one denotes by $Sect (Y,\int _{Y}\mm )_{c,f}$ the subcategory of objects
cofibrant and fibrant for the Reedy structure
(with the analogous notation for
$Z$) then one obtains the cartesian square
$$
\begin{array}{ccc}
Sect (Y,\int _{Y}\mm ) _{c,f}& \rightarrow & \prod _iSect (Y_i,\int _{Y_i}\mm )
_{c,f}
\\ \downarrow && \downarrow \\
Sect (Z,\int _{Z}\mm )  _{c,f}& \rightarrow &
\prod _i Sect (Z_i,\int _{Z_i}\mm )  _{c,f}.
\end{array}
$$
Here again, the horizontal morphisms respect weak equivalences.
It follows that the cartesian square
$$
\begin{array}{ccc}
Sect (Y,\int _{Y}\mm ) _{c,f}& \rightarrow & \prod _i \mm (y_i) ^{[0,1,2]}
_{c,f}\\ \downarrow && \downarrow \\
Sect (Z,\int _{Z}\mm )  _{c,f}& \rightarrow & \prod _i\mm (y_i) ^{[0,2]} _{c,f}.
\end{array}
$$

Let $N$ be one of the model categories in the preceding cartesian square
(either one of the factors of a product, or the product), and let
$N_{c,f}$ be its subcategory defined above. Under the hypothesis of
functorial factorisations, there exists a functor $\zeta : N\rightarrow
N_{c,f}$ and a natural transformation $t_u: \zeta
(u)\stackrel{\cong}{\rightarrow} u$. This functor respects weak
equivalences and the $t_u$ are weak equivalences. In particular, it
establishes an equivalence $L(N_{c,f})\cong L(N)$.



Let $\sigma , \tau \in Sect (Y, \int _Y\mm )_{c,f}$, and denote by
$\sigma |_Z$, $\tau |_Z$, $\sigma ^i_{012}$, $\tau ^i_{012}$, and $\sigma
^i_{02}$, $\tau ^i_{02}$ their images respectively in $Sect (Z,\int _Z\mm
|_Z)_{c,f}$, $\mm (y_i)^{[0,1,2]}_{c,f}$, and $\mm (y_i)^{[0,2]}_{c,f}$. We
claim that the square $(\ast )$
$$
\begin{array}{ccc}
L(Sect (Y, \int _Y\mm )_{c,f})_{1/}(\sigma , \tau )
&\rightarrow & \prod _i L(\mm (y_i)^{[0,1,2]}_{c,f})_{1/}(\sigma
^i_{012},\tau ^i
_{012})\\ \downarrow && \downarrow \\
L(Sect (Z, \int _Z\mm |_Z)_{c,f})_{1/}(\sigma |_Z, \tau |_Z)
& \rightarrow & \prod _i L(\mm (y_i)^{[0,2]}_{c,f})_{1/}(\sigma ^i_{02},\tau
^i_{02}) \end{array}
$$
is homotopy cartesian.

Suppose that we have shown that $(\ast )$ is homotopy cartesian,
and let us finish the proof of the theorem.
One has a homotopy cartesian square
$$
\begin{array}{ccc}
Sect (Y, \int _Y L(\mm ))_{1/}(\sigma , \tau )
&\rightarrow & \prod _i L(\mm (y_i))^{[0,1,2]}_{1/}(\sigma ^i_{012},\tau ^i
_{012})\\ \downarrow && \downarrow \\
Sect (Z, \int _Z L(\mm )|_Z)_{1/}(\sigma |_Z, \tau |_Z)
& \rightarrow & \prod _i L(\mm (y_i))^{[0,2]}_{1/}(\sigma ^i_{02},\tau
^i_{02}), \end{array}
$$
where for the direct products one first takes Kan replacements of the
simplicial sets (i.e. one takes the homotopy direct products). To justify
this, one makes a ``homotopy'' version for the $L(\mm )$ of the
discussion of diagrams over Reedy categories, using the notion
of homotopy adjoint of \S 8. See also the proof of
Theorem \ref{strictif1A}. We do not give further
details
on this point.

The morphism from the square $(\ast )$ into this one
induces an equivalence on the three terms
at the bottom and on the right. Indeed, for $Z$ one supposes
the theorem
known by recursion, and on the right it is a matter of diagrams
indexed by the direct categories $[0,1,2]$ and $[0,2]$---which we
suppose
already treated.
By Lemma
\ref{equivalence1}, the direct products in the square $(\ast )$ are
homotopy direct products.

The fact that the two squares are
homotopy cartesian then implies that the morphism
$$
L(Sect (Y, \int _Y\mm )_{c,f})_{1/}(\sigma , \tau )
\rightarrow Sect (Y, \int _Y L(\mm ))_{1/}(\sigma , \tau )
$$
is an equivalence, which proves the recursion step.

To finish the proof, we must: (1) justify the hypothesis that
the theorem is known for direct categories (of finite length,
this is what we use in the proof above); (2) show that the
square $(\ast )$ above is homotopy cartesian; and (3) justify the
passage to the limit over
the $k$ of the $F^kY$ for a Reedy category $Y$ with
unbounded degree function.

For (1), suppose now that one wants to prove the theorem for a
direct category $Y$ (of finite length). One proceeds in the same way
by recursion on the length of $Y$, and one copies the steps of the
proof above. The only change is that the categories
$Match(y_i)$ are empty and hence the terms $match (\sigma , y_i)$
do not appear. In particular, one may replace
the category of diagrams
$\mm (y_i)^{[0,1,2]}$ of length $3$ by a category of
diagrams $\mm (y_i)^{[0,1]}$ of length $2$; and the category of diagrams
$\mm (y_i)^{[0,1]}$ becomes just $\mm (y_i)$. This means that in the end
one is reduced to considering the case of $\mm (y_i)^{[0,1]}$,
in other words the case of diagrams indexed by $Y= I=[0,1]$ with values in
a fixed closed model category $M=\mm (y_i)$. This case was
treated by Lemma \ref{equivalence2}.

For (2), we shall again use the method of resolutions. First, one may
suppose that $\sigma$ is fibrant and cofibrant, as is $\tau$. One chooses
a simplicial resolution $\tau \rightarrow {\bf t}$ and one may
suppose that ${\bf t}$ is fibrant and cofibrant as a simplicial object,
i.e.
$$
{\bf t}\in Sect (Y,\int _Y\mm )^{\Delta ^o}.
$$
Note here that it is a matter of iterating twice the operation
which consists in ``taking the
Reedy structure'': first one took the Reedy structure on
$Sect (Y,\int _Y\mm )$ and then, with respect to this structure, one took
the Reedy structure on the simplicial diagrams inside it.

In particular, for every $p\in \Delta$, ${\bf t}(p)\in Sect (Y,\int _Y\mm )$
is both fibrant and cofibrant, hence
$$
{\bf t}\in [Sect (Y,\int _Y\mm )_{c,f}]^{\Delta ^o}.
$$
The images ${\bf t} |_Z$,  ${\bf t} ^i_{012}$, and  ${\bf t} ^i_{02}$,
which are simplicial objects in the other cmc's $N$ appearing above,
are in fact simplicial objects of the $N_{c,f}$.

One must fix a structure
of cmc on $\mm (y_i)^{[0,1,2]}$ and on $\mm (y_i)^{[0,2]}$. For this, one
will use a
Reedy structure on $[0,1,2]$ for which the object $1$ is of maximal
degree; thus the arrow $01$ is direct, and $12$ is inverse. Now there
is a choice to be made as regards the arrow $02$: we shall choose to say
that it is a direct arrow, which corresponds to saying $deg (0)<deg(2)$
(one must make the degree function ``tilt'' one way or the other).
The corresponding structure on $\mm (y_i)^{[0,2]}$ consists in saying
again that the
arrow $02$ is direct. With this choice, an object of $\mm (y_i)^{[0,2]}$ is
fibrant if and only if its components are fibrant, while an object is
cofibrant if and only if its components are fibrant % source prints "fibrantes"; kept as printed
and the arrow $02$
is a cofibration. An object of $\mm (y_i)^{[0,1,2]}$ is fibrant if and
only if
its components as well as the arrow $12$ are; and an object is cofibrant if
and only if its components as well as the arrows $01$ and $02$ are.

One still has that $\tau ^i_{012}$ and ${\bf t}^i_{012}$ are fibrant objects
(likewise for $\tau ^i_{02}$ and ${\bf t}^i_{02}$). By contrast, $\sigma
^i_{012}$ and $\sigma ^i_{02}$ are no longer cofibrant since one cannot
guarantee that the arrow $\sigma ^i_0\rightarrow \sigma ^i_2$ be a
cofibration.
To arrange this one chooses a cofibrant replacement
$\tilde{\sigma}^i_{012}\rightarrow  \sigma ^i_{012}$ in the following
way: one
chooses a diagram
$$
\begin{array}{ccc}
\sigma ^i_1 & \rightarrow & \tilde{\sigma}^i_2\\
\uparrow & \nearrow & \downarrow \\
\sigma ^i_0 & \rightarrow & \sigma ^i_2
\end{array}
$$
with a trivial fibration on the right, and
a cofibration
on the diagonal. The left-hand morphism remains a cofibration and we say
nothing about the top morphism. The composite $\sigma ^i_1\rightarrow \sigma ^i_2$
should be the one we already have. One sets
$$
\tilde{\sigma}^i_{012}:= [\sigma ^i_0\rightarrow \sigma ^i_1 \rightarrow
\tilde{\sigma}^i_2].
$$
The square of simplicial sets
$$
\begin{array}{ccc}
\mm (y_i)^{[0,1,2]}(\sigma ^i_{012},{\bf t} ^i
_{012})
&\rightarrow & \mm (y_i)^{[0,1,2]}(\tilde{\sigma} ^i_{012},{\bf t} ^i
_{012})\\ \downarrow && \downarrow \\
\mm (y_i)^{[0,2]}_{1/}(\sigma ^i_{02},{\bf t}
^i_{02})
& \rightarrow & \mm (y_i)^{[0,2]}_{1/}(\tilde{\sigma} ^i_{02},{\bf t}
^i_{02})
\end{array}
$$
is cartesian. Likewise the square
$$
\begin{array}{ccc}
Sect (Y, \int _Y\mm )_{1/}(\sigma , {\bf t} )
&\rightarrow & \prod _i \mm (y)^{[0,1,2]}(\sigma ^i_{012},{\bf t} ^i
_{012})\\ \downarrow && \downarrow \\
Sect (Z, \int _Z\mm |_Z)_{1/}(\sigma |_Z, {\bf t} |_Z)
& \rightarrow & \prod _i \mm (y)^{[0,2]}_{1/}(\sigma ^i_{02},{\bf t}
^i_{02}) \end{array}
$$
is cartesian and by composing with the product over $i$ of the preceding
squares
one obtains the cartesian square $(\ast \ast )$ of simplicial sets
$$
\begin{array}{ccc}
Sect (Y, \int _Y\mm )_{1/}(\sigma , {\bf t} )
&\rightarrow & \prod _i \mm (y_i)^{[0,1,2]}(\tilde{\sigma} ^i_{012},{\bf t} ^i
_{012})\\ \downarrow && \downarrow \\
Sect (Z, \int _Z\mm |_Z)_{1/}(\sigma |_Z, {\bf t} |_Z)
& \rightarrow & \prod _i \mm (y_i)^{[0,2]}_{1/}(\tilde{\sigma} ^i_{02},{\bf t}
^i_{02}).
\end{array}
$$
This square is equivalent to the square $(\ast )$
(to find the morphism between squares, note in general that
if $(C,W)$ is a category with a subcategory of ``equivalences'' and if
$x,y$
are objects of it, with an augmented simplicial object $y\rightarrow {\bf y}$
such that for every $p$ the morphism $y\rightarrow {\bf y}(p)$ lies in $W$,
then one obtains a natural morphism of simplicial sets
$C_{1/}(x,{\bf y})\rightarrow L^H(C,W)_{1/}(x,y)$ where $L^H$ is the
localisation by hammocks \cite{DwyerKan2} \cite{DwyerKan3}; one applies this
to the
categories $C=N_{c,f}$ which appear in $(\ast )$).

Hence for (2) it suffices to see that $(\ast \ast )$ is
homotopy cartesian,
so it suffices to see that for every $i$ the morphism
$$
\mm (y_i)^{[0,1,2]}(\tilde{\sigma} ^i_{012},{\bf t} ^i
_{012})\rightarrow
\mm (y_i)^{[0,2]}_{1/}(\tilde{\sigma} ^i_{02},{\bf t}
^i_{02})
$$
is a Kan fibration.

For this
we shall simplify the notation: one takes a cmc $M$ and two diagrams
$x,y\in M^{[0,1,2]}$ plus a fibrant simplicial resolution $y\rightarrow
{\bf y}$. For this one equips (as before) $M^{[0,1,2]}$ with a cmc structure
such that the fibrant objects are the diagrams of fibrant objects with the
second morphism $12$ fibrant, and the cofibrant objects are the diagrams
of cofibrant objects with the morphisms $01$ and $02$ cofibrant. Likewise one
equips $M^{[0,2]}$ with the structure for which the cofibrant objects are
the diagrams of cofibrant objects with a cofibration in $02$.

Let us denote by $i$ the morphism $ [0,2]\hookrightarrow [0,1,2]$; one has the
functor
$$
i^{\ast}: M^{[0,1,2]}\rightarrow M^{[0,2]}
$$
and its left adjoint $i_!$. This adjoint has the following concrete
description. If $u= (u_0\rightarrow u_2)\in M^{[0,2]}$ then
$$
i_!(u)= [u_0 \stackrel{=}{\rightarrow} u_0 \rightarrow u_2]\in M^{[0,1,2]}.
$$
If $x$ is a cofibrant object of $M^{[0,1,2]}$ then
the adjunction morphism $i_!i^{\ast}x\rightarrow x$ is
$$
[x_0 \stackrel{=}{\rightarrow} x_0 \rightarrow x_2]
\rightarrow [x_0\rightarrow x_1\rightarrow x_2]
$$
which is a cofibration.

One has
$$
M^{[0,2]}_{1/}(i^{\ast} x, i^{\ast} {\bf y})=
M^{[0,1,2]}_{1/}(i_!i^{\ast} x, {\bf y}).
$$
The morphism in question is therefore the morphism induced by
$i_!i^{\ast}x\rightarrow x$,
$$
M^{[0,1,2]}_{1/}(x, {\bf y})
\rightarrow
M^{[0,1,2]}_{1/}(i_!i^{\ast} x, {\bf y}).
$$
If $x$ is cofibrant then the cofibration $i_!i^{\ast}x\rightarrow x$
induces here a Kan fibration of simplicial sets,
which shows that $(\ast \ast )$ is homotopy cartesian,
whence the same thing
for $(\ast )$, which gives (2).



By the way, the argument above shows that if $\sigma$ and $\tau$ (resp. ${\bf
t}$) are cofibrant and fibrant sections (resp.
is a fibrant
and cofibrant simplicial
resolution) then the morphism
$$
Sect (Y, \int _Y\mm )_{1/}(\sigma , {\bf t} )\rightarrow
Sect (Z, \int _Z\mm |_Z)_{1/}(\sigma , {\bf t} )
$$
is a Kan fibration.

For (3) one now fixes a Reedy category $Y$ with a degree
function not necessarily bounded. One fixes a fibrant replacement
$$
\int '_YL(\mm )\rightarrow Y
$$
which determines fibrant replacements by restriction
$$
\int '_{F^kY} L(\mm )|_{F^kY}\rightarrow F^kY.
$$
Under these conditions, $Sect (Y, \int '_YL(\mm ))$ is the limit of the
sequence
$$
\ldots \rightarrow Sect (F^kY, \int '_{F^kY}L(\mm )|_{F^kY})
\rightarrow Sect (F^{k-1}Y, \int '_{F^{k-1}Y}L(\mm
)|_{F^{k-1}Y})\rightarrow \ldots ,
$$
and the transition morphisms in this sequence are fibrations
of Segal $1$-categories. On the other hand, let us fix sections $\sigma$ and
$\tau$ in $Sect (Y, \int _Y\mm )$, with $\sigma$ cofibrant, and fix a
fibrant simplicial resolution $\tau \rightarrow {\bf t}$. One has that
$$
Sect (Y, \int _Y\mm )_{1/}(\sigma , {\bf t})
$$
is the limit of the sequence of simplicial sets
$$
\ldots \rightarrow Sect (F^kY, \int _{F^kY}\mm |_{F^kY})_{1/}(\sigma
|_{F^kY}, {\bf t}|_{F^kY})
\rightarrow \ldots .
$$
By the argument given for (2), the transition morphisms
of this sequence are
Kan fibrations. This sequence is, term by term, equivalent
(by the same argument) to the sequence
$$
\ldots \rightarrow Sect (F^kY, \int '_{F^kY}L(\mm )|_{F^kY})_{1/}(\sigma
|_{F^kY}, \tau |_{F^kY}
) \rightarrow \ldots ,
$$
and these equivalences being functorial one obtains an equivalence
of diagrams between the two sequences (in fact, to really have a morphism
here it would be appropriate to use the localisation by hammocks $L^H$
instead of $L$). The
fact that the transition morphisms in the two cases are Kan
fibrations implies that the morphism on the limits is an equivalence
$$
Sect (Y, \int _Y\mm )_{1/}(\sigma , {\bf t})\stackrel{\cong}{\rightarrow}
Sect (Y, \int '_YL(\mm ))_{1/}(\sigma , \tau ).
$$
Now
$$
Sect (Y, \int _Y\mm )_{1/}(\sigma , {\bf t}) \cong
L(Sect (Y, \int _Y\mm ))_{1/}(\sigma , \tau )
$$
whence part (3).

This completes the proof of Theorem \ref{uneequivalence}.
\eop
